\section{Introduction}

The Supreme Court's decision in \emph{Rucho v. Common Cause} (2019) declared partisan gerrymandering claims non-justiciable, removing federal judicial oversight of redistricting practices that Chief Justice Roberts acknowledged ``incompatible with democratic principles.'' In \emph{Rucho}'s wake, reformers have increasingly turned to algorithmic redistricting as a potential solution: if districts are drawn by algorithms that cannot access partisan data, the resulting maps cannot constitute intentional partisan gerrymanders~\cite{cho2016measuring,duchin2018geometry}.

However, this ``impossibility defense'' faces a fundamental challenge: geography itself encodes partisan patterns. Democratic voters concentrate in dense urban cores while Republican voters disperse across suburbs and rural areas~\cite{rodden2019cities,chen2013unintentional}. An algorithm optimizing for geographic compactness will inevitably pack urban Democrats into fewer districts while splitting rural Republicans across more districts, producing partisan asymmetry without partisan intent. Critics argue that algorithmic redistricting merely replaces human manipulation with geographic determinism, offering procedural objectivity but not substantive partisan fairness~\cite{stephanopoulos2015partisan}.

This paper addresses a critical empirical question: \emph{How much partisan bias do algorithmic methods produce, and how does this compare to enacted plans?} While prior work has demonstrated that algorithmic redistricting can satisfy constitutional requirements~\cite{deluca2026recursive} and exceed Voting Rights Act compliance~\cite{deluca2026vra}, the partisan fairness implications remain underexplored. Do neutral algorithms reduce partisan bias relative to human-drawn maps, or do they simply reflect—and potentially entrench—existing geographic asymmetries?

We answer this question through the first national-scale efficiency gap analysis of algorithmic redistricting. The efficiency gap, developed by Stephanopoulos and McGhee~\cite{stephanopoulos2015partisan}, measures partisan bias by calculating ``wasted votes'': ballots cast for losing candidates or surplus votes beyond the threshold needed to win. By comparing wasted votes between parties, the efficiency gap quantifies whether a redistricting plan systematically advantages one party over another. We compute efficiency gaps for algorithmic plans (generated via recursive bisection that cannot access partisan data) and enacted plans across all 50 states for the 2016-2020 election cycle, enabling direct comparison of partisan fairness.

Our analysis reveals three key findings. First, algorithmic redistricting exhibits substantially lower partisan bias than enacted plans: mean efficiency gap of $-3.2\%$ (slight Democratic advantage) versus $+5.1\%$ (Republican advantage), a difference of 8.3 percentage points representing 62\% bias reduction. Second, regional variation is pronounced: Rust Belt states (Pennsylvania, Wisconsin, Michigan) show the largest disparities, with enacted plans averaging $+7.2\%$ efficiency gaps favoring Republicans while algorithmic plans average $-2.8\%$ favoring Democrats. Sunbelt states show smaller but consistent patterns. Third, the persistent negative efficiency gap in algorithmic plans—even when optimizing solely for geographic compactness—establishes an empirical lower bound on partisan asymmetry achievable through neutral redistricting. This $-3.2\%$ baseline reflects unavoidable Democratic urban concentration and cannot be eliminated without either relaxing compactness requirements or explicitly targeting partisan balance.

These findings have important implications for redistricting reform and legal doctrine. For reformers, our results demonstrate that algorithmic methods provide meaningful improvements in partisan fairness while maintaining other redistricting values (compactness, contiguity, population equality). For courts, the efficiency gap benchmarks provide quantitative standards for evaluating whether enacted plans exhibit bias beyond geographic baselines. If a state's enacted plan shows $+7\%$ efficiency gap while algorithmic plans for that state average $-3\%$, the 10-percentage-point difference suggests human manipulation rather than geographic inevitability. Finally, our results clarify the limits of algorithmic objectivity: neutral algorithms reduce but cannot eliminate partisan effects, and reformers seeking proportional representation may need to consider alternatives to single-member districts~\cite{mcghee2020fairness}.

The paper proceeds as follows. Section~2 reviews the efficiency gap literature and prior algorithmic redistricting research. Section~3 describes our methodology for generating algorithmic plans and computing efficiency gaps. Section~4 presents national and regional results comparing algorithmic versus enacted plans. Section~5 discusses mechanisms driving geographic bias and implications for reform. Section~6 concludes by considering limitations and future directions for quantitative partisan fairness analysis.
